% 本文件由 LaTeX 排版 Agent 根据有效 translation.json 自动生成。
% author=Tertullian; work_id=0321; title=\论圣洗

\chapter{\WorkTitle{论圣洗}}

% structure: p0001 source=h1 target=omitted reason=page-title-already-rendered-by-parent

% structure: p0002 source=h2 target=section reason=relative-heading-rank; base=h2

\section{第一章 引言 本论著的起源}

我们的水圣事是有福的，因为它洗去我们早期盲目之罪，我们就得释放，并\Emph{被接纳}进入永生！关于此事的一篇论著并非多余；它不仅教导那些刚在信德中成形的人，也教导那些满足于简单相信、没有充分查考传统根据的人，他们因无知而怀有一种未经检验\Emph{却}似可信的信德。结果，近来在这一带活动的一位加音异端的毒蛇，以她最具毒害的教义掳走了许多人，其首要目标就是摧毁圣洗。这完全合乎本性；因为毒蛇、角蝰和蛇怪本身通常喜欢干旱无水之地。但我们，小鱼，效法我们的\OriginalTerm{鱼}{\GreekText{ΙΧΘΥΣ}}耶稣基督，在水中出生，除了永久留在水中，别无安全之道；因此，那个最可怕的受造物，她甚至没有权利教导健全的道理，却深知如何通过将小鱼带离水来杀死它们！\par

% structure: p0004 source=h2 target=section reason=relative-heading-rank; base=h2

\section{第二章 天主工作方法的极度单纯——属肉性心思的绊脚石}

然而，乖谬的力量何其大，竟能\Emph{如此}动摇信德或完全阻止人接受信德，以致它正是在构成\Emph{信德}的原则上攻击信德！没有什么比可见于行为的天主工程之单纯，与其效果中所应许的宏伟相比，更能使人心顽硬；正因如此，基于这样大的单纯——没有排场，没有任何特别的准备，最后，没有花费——一个人浸入水中，念几句简短的话时被洒水，然后起来，并未干净多少（或完全没有），随之而来的永生之获得便更被视为不可思议。若不然，我就是骗子：相反，正是由于环境、准备和花费，偶像的庆典或神秘仪式才建立起它们的信誉和权威。哦，可怜的不信！竟完全否认天主自己的属性：单纯与全能！那么，难道死亡被沐浴洗去不也是奇妙的吗？但正因为奇妙，它才更可信，如果奇妙正是它\Emph{不}被信的理由。因为天主的工程，若不超越一切奇妙，又该具有什么性质呢？我们也感到惊奇，但这是\Emph{因为}我们相信。而不信也同样惊奇，但它并不相信：它惊奇于单纯\Emph{行为}，视其为虚妄；惊奇于宏大\Emph{效果}，视其为不可能。即使你认为这足以应对每一点，但已有天主的先行宣告：\BibleQuote{天主拣选了世上愚拙的事，以挫败智慧；}和\BibleQuote{在人看来极难的事，在天主却是容易的。}因为如果天主是明智和全能的——即使那些忽视他的人也不否认——那么他有充分理由将他自身运作的物质原因置于智慧与力量的对立面，即愚拙与不可能之中；因为每一种德能都从其被激发的事物中获得原因。\par

% structure: p0006 source=h2 target=section reason=relative-heading-rank; base=h2

\section{第三章 水被选为天主运作的器皿及其原因——它在创造中居首位}

既然铭记这项声明如同一个决定性的指令，我们仍要\Emph{探究}这个\Emph{问题}：\Quote{藉水重新塑造是何等\Emph{愚昧}和\Emph{不可能}？}请问，这种物质材料凭什么配得上如此崇高的职分？我想，有必要考察液体元素的权威。然而，这种权威从一开始就大量存在。因为\Emph{水}是在世界一切装饰之前，与天主一同处于未成形状态的事物之一。\Quote{在起初，}\Emph{圣经}说，\Quote{天主创造了天地。但地是混沌空虚，深渊上一片黑暗，天主的神在水面上运行。}人啊，你首先应敬重的是水的\Emph{古老}，因为它们的实体是古老的；其次是它们的\Emph{尊贵}，因为它们是天主圣神的居所，无疑比当时其他一切存在的元素更\Emph{中悦于祂}。因为那时黑暗是彻底的，无形无状，没有星辰的装饰；深渊是阴暗的；地是未装备的；天是未加工的：唯独水——始终是一种完美、喜乐、简单的物质材料，本身纯洁——为天主提供了适宜的载体。\Emph{事实是}，水在某种程度上是规范的力量，天主藉此制定了世界的后续安排？因为祂藉着\BibleQuote{分开水}使穹苍悬于中间；藉着\BibleQuote{分出水}使\BibleQuote{旱地}出现。此后，世界通过其元素被安置，当居民被给予时，\BibleQuote{水}是第一个接受命令\BibleQuote{生出活物}的。水首先产生了有生命之物，这样在圣洗中，水懂得赋予生命就不足为奇了。难道连人的形成不也是借助水的帮助完成的吗？合适的材料在于\Emph{土}，但除非湿润多汁，否则不适于这个目的；而\BibleQuote{水}在第四天之前被分离到自己的地方，用它们残留的水分调制成黏土般的稠度。从那时起，如果我继续全面或更详细地列举我能引述的关于这种元素\Emph{权威}的证据，以显示其力量或恩宠何其大，它给世界提供了多少巧妙的发明、多少功能、多少有用的工具，恐怕我似乎是在搜集水的赞美而非圣洗的理由；尽管我\Emph{由此}会更充分教导，天主使祂在其所有产品和作为中安排的物质材料，也在祂特有的圣事中服从祂；那种统治地上生命的\Emph{物质材料}也作用于天上的生命。\par

% structure: p0008 source=h2 target=section reason=relative-heading-rank; base=h2

\section{第四章　圣神在水上运行的原始典型与圣洗圣事。水这普遍元素因而成为圣化的管道。外在标记与内在恩宠的相似。}

但这样就足够在一开始指出那些要点，其中甚至承认圣洗的首要原则——由那为圣洗预像而\Emph{采取}的姿态预先注定的——即从一开始就运行在水上的天主圣神，要继续停留在受洗者的水上。当然，圣物运行在圣物之上；否则，从运行者那里，被运行者借得圣洁，因为在每种情况下，底层的物质材料必然会吸收其上覆物的性质，尤其是物质的吸收精神的，通过其物质的微妙性，既渗透又渗入。如此，被圣者圣化的水的本性，自身也获得了圣化的能力。没有人可以说：\Quote{那么请问，我们为何用最初存在的水受洗？}当然不是用那些水，除非\OriginalTerm{属}{genus}确实是一个，但\OriginalTerm{种}{species}却很多。然而，属于\OriginalTerm{属}{genus}的特性也出现在\OriginalTerm{种}{species}中。因此，一个人在海洋或池塘、溪流或泉水、湖泊或水槽中受洗，并无区别；若翰在约旦河施洗的那些人和伯多禄在台伯河施洗的那些人之间也没有区别，除非连斐理伯在路上用偶然的水为太监施洗，他从那水所得救恩并不比别人多或少。\BibleRef{宗 8:26}因此，所有水，凭借其起源的原始特权，在呼求天主后，都获得圣化的圣事性能力；因为圣神立即从天而降，运行在水上，从自身圣化它们；如此被圣化后，它们同时汲取了圣化的能力。尽管可以承认这种相似性适合简单的行为：既然我们被罪玷污，如同被污秽，我们应藉水洗净这些污点。但罪并不显现在我们的\Emph{肉体}中（因为没有人会在皮肤上带着拜偶像、奸淫或欺诈的痕迹），所以那类人在\Emph{灵魂}中是污秽的，因为灵魂是罪的根源；灵魂是主人，肉体是仆人。然而，两者互相分担罪责：灵魂基于命令，肉体基于服从。因此，在水通过天使的中介获得药力后，灵魂在身体上被水洗净，肉体在同一水中得到精神上的洁净。\par

% structure: p0010 source=h2 target=section reason=relative-heading-rank; base=h2

\section{第五章　外教人对水的运用。贝特匝达池天使的预像。}

\Quote{唉，但那些对属灵能力毫无认识的异邦人，却将同样的功效归诸他们的偶像，认为水具有这种性能。}（他们确实如此）但他们以无丈夫之水自欺。因为洗涤是他们加入某些神秘仪式——比如那臭名昭著的伊西斯或密特拉——的途径。他们也藉洗涤来尊崇神祇本身。此外，他们用运水洒水的方式，为乡间别墅、房屋、庙宇乃至整座城市举行涤罪；无论如何，在阿波罗节和厄琉西斯节上，他们受洗；并自以为这样做的效果就是他们的重生，以及因背誓而应受的惩罚得到赦免。同样，在古代，凡以杀人玷污自己的人，常去寻找净化之水。因此，如果仅仅是水的本性——即它是合适的洗涤材料——就使人们自欺地相信净化预兆的话，那么，水通过那构成其全部本性的天主的权威，岂不更真实地发挥这种作用？如果人们认为水因宗教而具有医治功效，那么还有什么宗教比生活天主的宗教更有效呢？既然认清了这一点，我们在此也看到魔鬼与天主的事物争竞的热忱，我们发现他也正在他的\Emph{臣属}身上施行洗礼。何等相似！不洁者净化！毁灭者释放！被定罪者赦免！他真要借洗去他自己所激励的罪恶来毁灭自己的工程吗？这些话是作为对不信者所做的证言；如果他们不信天主的事物，却相信那在仇敌天主者身上的虚假仿效。难道没有其他例子吗？在那时，没有圣事，不洁的神灵也在水上盘旋，仿效起初天主圣神在水上的盘旋。请看一切阴暗的泉源、一切人迹罕至的溪流、浴池中的池塘、\Emph{私人}房屋的水渠，或那些据说具有\Quote{摄魂}能力的水池和水井，这能力来自那行恶神灵。凡被水淹死或受其影响而发疯或恐惧的人，他们称之为被水仙精所迷，或\Quote{\OriginalTerm{水狂}{lymphatic}}，或\Quote{\OriginalTerm{惧水}{hydro-phobic}}。我们为什么举这些例子？免得有人难以\Emph{相信}：天主的圣天使会临于水上，使之柔和以助人得救；而恶天使却常常亵渎地利用同一元素来毁灭人。如果以为天使临于水是一种新奇的事，那么已有未来的预兆为先例。天使曾以其干预搅动贝特匝达的水池。患病的人常等待他；因为谁先下去，洗后便痊愈了。这肉体的治愈形象预示了精神的治愈，按照规则：属血肉的事常先于属神的事，作为预象。如此，当天主的恩宠在人间进展到更高程度时\BibleRef{若 1:16}，水与天使获得了\Emph{效力}的增加。那些曾经医治身体缺陷的，如今治愈灵魂；那些曾赐予暂时救恩的，如今赐予永恒；那些一年只释放一次的，如今每日整批拯救万民，罪恶的洗涤消除了死亡。罪债既除，惩罚当然也除去了。于是，人将被恢复给天主，成为祂的\Quote{相似}，而先前他曾被\Emph{塑造}为天主的\Quote{肖像}；（\Quote{肖像}被视为在于他的\Emph{外形}，\Quote{相似}在于他的\Emph{永恒}；）因为他再次领受了天主的圣神，那最初从祂的\Emph{气息}所领受，却因罪而失落的圣神。\par

% structure: p0012 source=h2 target=section reason=relative-heading-rank; base=h2

\section{第六章 天使——圣神的先驱；洗礼规式中所含的意义}

并非\Emph{在}水中我们获得圣神；而是在水中，在天使的（见证）下，我们被洗净，并为圣神而\Emph{预备}。在此也有一个预象先行：因为这样，若翰事先作了主的先驱，\Quote{预备祂的道路}\BibleRef{路 1:76}。同样，作为洗礼见证的天使，也为那将要降临于我们的圣神\Quote{修直路径}，藉着洗去罪恶，而这罪恶正是因信德——在圣父、圣子和圣神的名下所印证的——所获得的。因为如果\Quote{\Emph{凭三个见证}的口，每句话都应成立}——而我们藉着祝福，拥有同样的（三位）作为我们信德的见证，同时也是我们救恩的保证——那么，神圣的名字的数量岂不更足以担保我们的希望吗！此外，在\Quote{三个见证}之下既已作出了信德的见证和救恩的许诺后，必当提及教会；因为凡有三者（即圣父、圣子和圣神）之处，就有教会，那是三者的身体。\par

% structure: p0014 source=h2 target=section reason=relative-heading-rank; base=h2

\section{第七章 论傅油}

之后，我们从池中出来，便用祝福过的油彻底傅抹——（这源于）古代的礼规，在那时，人们进入司铎职时，惯用角中的油傅抹，自亚郎被梅瑟傅油以来便是如此。因此亚郎被称为\Quote{基督}，源自\Quote{基督圣油}，即\Quote{傅油}；这傅油，当成为属灵的时，便为主提供了恰当的名字，因为祂被圣父以圣神\Quote{傅油}；正如\WorkTitle{宗徒大事录}所\Emph{记载}：\Quote{他们确实在这城里聚集，反对祢所傅油的圣子。}同样，在\Emph{我们}身上，傅油按肉体进行（即在身体上），但在灵性上获益；正如洗礼本身的\Emph{行为}也是属肉体的，即我们浸入水中，\Emph{而}其\Emph{效果}是属灵的，即我们获得罪恶的赦免。\par

% structure: p0016 source=h2 target=section reason=relative-heading-rank; base=h2

\section{第八章  按手礼。洪水与鸽子的预象}

接下来，按手在我们身上，藉着祝福呼求并邀请圣神。难道人的聪明能将一个灵召入水中，藉着从上按手，使它们与另一个如此清澈的灵结合为一个身体吗？难道天主就不能在祂自己的器皿上，藉着\Quote{圣手}产生一种崇高的属灵旋律吗？但这一点，如同前者，也源自古老的圣事礼仪，其中雅各伯为他若瑟所生的孙子厄弗辣因和默纳协祝福；他将手按在他们身上并交换，且如此斜斜地交叉，以致于藉此描绘基督，甚至预示了将来在基督内的祝福。然后，在我们洁净并蒙福的身体上，那至圣的圣神甘愿从圣父降临。在圣洗的水上，如同认出祂原始的居所，祂安息；祂曾以\Quote{鸽子的形状}降临在主身上，为藉着纯朴与无罪的受造物（象征）宣告圣神的性质，因为鸽子在其身体结构中没有苦胆。因此祂说：\Quote{你们要纯朴如同鸽子。}这一点也并非没有先前预象的佐证。因为正如洪水之后——这洪水除去了旧日的邪恶，可以说是世界的洗礼——一只\Emph{鸽子}从方舟中被派出去宣告天主愤怒的平息，她带着橄榄枝归来，这标记甚至在万民中也是\Emph{和平}的前兆；同样，按照同一属天效果的法则，圣神的\Emph{鸽子}从天上——那里有教会，预象的方舟——飞出，临到从洗礼池中升起、除去了旧罪的大地（即我们的肉体），给我们带来天主的和平。但世界重新陷于罪恶；在这一点上，洗礼很难与洪水相比。因此它注定要受火刑；正如人在领洗后重新犯罪也是如此：故此，这也应被接纳为劝诫我们的一个标记。\par

% structure: p0018 source=h2 target=section reason=relative-heading-rank; base=h2

\section{第九章  红海与磐石出水的预象}

自然有多少理由，恩宠有多少特权，纪律有多少仪式、预象、预备、祈祷，确立了水的圣洁？首先，当人民无条件获得自由，藉着穿越\Emph{水}逃脱埃及王的暴力时，正是\Emph{水}淹没了那王及其全军。\BibleRef{出 14:27}还有什么预象比这在洗礼圣事中更明显地应验？万民藉着\Emph{水}从世界获得自由，即：他们完全撇下那旧暴君魔鬼，使其淹没在\Emph{水}中。再者，\Emph{水}藉着梅瑟的树木从\Quote{苦}的缺陷恢复其本有的\Quote{甜}的恩宠。那树木就是基督，祂将自己恢复曾被毒化、苦涩的本性之\Emph{泉脉}成为全然救恩的\Emph{洗礼之水}。这就是那从\Quote{伴随的磐石}不断为百姓流出的\Emph{水}；因为如果基督是\Quote{磐石}，那么无疑我们看见洗礼因着基督里的\Emph{水}而得祝福。在天主及祂的基督眼中，为了坚定洗礼，\Emph{水}的恩宠何其伟大！基督从未没有\Emph{水}：即，祂自己在\Emph{水}中受洗；\BibleRef{玛 3:13}在加纳婚宴受邀时，以\Emph{水}开始祂权能的初步显示；\BibleRef{若 2:1}在讲论时，祂邀请口渴的人到祂自己的永恒\Emph{水}那里；\BibleRef{若 7:37}当教导爱德时，祂赞许在慈善工作中给穷人一杯\Emph{水}；\BibleRef{玛 10:42}祂在\Emph{井}旁恢复体力；\BibleRef{若 4:6}在\Emph{水}上行走；\BibleRef{玛 14:25}甘愿渡过\Emph{海}；\BibleRef{谷 4:36}给门徒\Emph{水}。\BibleRef{若 13:1}直到受难，洗礼的见证持续：当祂被交出钉十字架时，\Emph{水}介入，比拉多洗手为证；当祂被刺伤时，从祂肋旁涌出\Emph{水}，士兵的枪为证！\par

% structure: p0020 source=h2 target=section reason=relative-heading-rank; base=h2

\section{第十章  若翰的洗礼}

我们已就我们有限的能力说到构成洗礼圣洁基础的\Emph{一般原则}。现在，同样尽我们所能，我将继续论述其特性的其余部分，涉及某些次要问题。\par

若翰所宣布的圣洗，甚至在那时，就由主亲自向法利塞人提出一个问题：那圣洗是属天的，还是真正属地的？关于此，他们无法给出一致的答案，因为他们不明白，因为他们不信。但\Emph{我们}，以同样贫乏的理解力和信德，\Emph{能够}确定那圣洗确实是\Emph{属神的}（然而就命令而言，并非也就在功效方面，因为我们读到若翰是\Emph{为主所派遣}来执行此职责的），但其本质是\Emph{属人的}：因为它并不传递任何属天的事物，而是预先服务于属天的事物；即，它是设立在\Quote{悔改}之上，这悔改是在人的能力之内的。事实上，律法经师和法利塞人，他们不愿\Quote{相信}，也未曾\Quote{悔改}。但如果悔改是人的事，那么它的圣洗必然也具有同样的性质；否则，如果它是属天的，它就会同时赐下圣神和罪的赦免。但除了天主以外，没有人能赦免罪或自由地赐予圣神。连主自己也说过，圣神不会在其他条件下降临，除非祂先升到父那里去。\BibleRef{若 16:6}主尚未赐予的，仆人当然不能提供。因此，在\WorkTitle{宗徒大事录}中，我们发现那些领受\Quote{若翰的圣洗}的人没有领受圣神，他们甚至连听闻都没有听过祂。那么，那不能提供属天的（恩赐）的，就不是属天的事物；然而，若翰身上真正属天的——预言之神——在全部圣神转移给主之后，就完全消失了，以至于他立刻派人去询问，他所亲自宣讲的、他所指出那来就他的那一位，是否就是\Quote{祂}。因此，\Quote{悔改的圣洗}\BibleRef{宗 19:4}被当作一位候选人，为即将在基督内到来的赦免和圣化而预备：因为若翰曾宣讲\Quote{为赦免罪的圣洗}\BibleRef{谷 1:4}，这宣认是针对\Emph{将来的}赦免；如果悔改在先、赦免在后是真的（事实就是如此），而这就是\Quote{预备道路}\BibleRef{路 1:76}。但那位\Quote{预备}的，并不亲自\Quote{成全}，而是为另一位得以成全而预备。若翰亲自承认，属天的事不是他的，而是基督的，他说：\Quote{那从地上来的，讲论地上的事；那从天上来的，在万有之上；}又说，他\Quote{只是以悔改施圣洗，但有一位即将到来，祂要以圣神和火施圣洗；}——当然，因为真实而稳固的信仰是以\Emph{水}受圣洗，以致救恩；虚假而软弱的信仰是以\Emph{火}受圣洗，以致审判。\par

% structure: p0023 source=h2 target=section reason=relative-heading-rank; base=h2

\section{第十一章 回答反对意见：\Quote{主没有施圣洗。}}

但是，有人说：\Quote{主来了，却没有施圣洗；因为我们读到，\InnerQuote{然而祂不施圣洗，而是祂的门徒施圣洗！}}\BibleRef{若 4:2} 仿佛若翰曾预言祂会亲手施圣洗！当然，他的话不应如此理解，而是按照普通方式说的；例如，我们说：\Quote{皇帝颁布了法令，}或\Quote{总督杖责了他。}请问，皇帝亲自颁布了吗？总督亲自杖责了吗？一个人由他的手下做事，总说是他做的。因此，\Quote{祂将给你们施圣洗}必须理解为指\Quote{借着祂}或\Quote{归于祂}，\Quote{你们将受圣洗。}但不要让\Quote{祂自己没有施圣洗}这一事实困扰任何人。因为祂要施圣洗归于谁呢？归于悔改？那么，你们把祂的前驱置于何用？归于罪的赦免，祂常用一句话赐予？归于祂自己，祂因谦卑而隐藏了？归于圣神，祂尚未从父那里降临？归于教会，祂的宗徒尚未建立？因此，祂的门徒正是用同一个\Quote{若翰的圣洗}来施圣洗，作为仆人，若翰先前也作为前驱用这个圣洗施圣洗。不要以为是用别的，因为除了后来基督的圣洗之外，没有别的；那时祂的门徒当然不能给予，因为主的荣耀尚未完全获得，圣洗池的功效尚未通过苦难和复活而确立；因为我们的死亡除非通过主的苦难，否则不能解脱，我们的生命除非通过祂的复活，否则不能恢复。\par

% structure: p0025 source=h2 target=section reason=relative-heading-rank; base=h2

\section{第十二章 论圣洗对救恩的必要性}

然而，当这一规定被提出：\Quote{没有圣洗，谁都得不到救恩}（主要基于主自己的宣告，祂说：\Quote{除非由水而生，他没有生命}），立刻有些人心生疑虑，甚至可以说是冒昧的质疑：\Quote{按此规定，宗徒们如何得救？他们中除了保禄，我们并未发现他们在主内受了圣洗？既然保禄是他们中唯一穿上基督圣洗之衣的，要么所有缺乏基督之水的人的得救危险被预先断定，以维持该规定；要么若连未受圣洗者也注定得救，则该规定被废除。}我听到过——主是我的见证——类似的疑虑：免得有人以为我如此放肆，无缘无故地纵笔构思出使他人心生困惑的想法。\par

如今，我尽我所能，要回答那些声称\Quote{宗徒们未受圣洗}的人。因为若他们领受了若翰的人间圣洗，并渴望主的圣洗，\Emph{那么}，既然主亲自规定圣洗是\Emph{一}（祂对渴望彻底沐浴的伯多禄说：\Quote{一次沐浴过的，无需\Emph{再洗}第二次。}\BibleRef{若 13:9}——当然，祂根本不会对一位\Emph{未}受圣洗的人这样说），我们在此就有一个明显的证据，反驳那些为毁灭水的圣事而连若翰的圣洗也剥夺宗徒的人。难道可以相信，\Quote{主的道}，即若翰的圣洗，那时尚未在那些注定要在全世界\Emph{开辟}主道的人身上\Quote{预备}好？主自己，虽然祂无需\Quote{悔改}，却受了圣洗：难道罪人不需要圣洗吗？至于\Quote{其他人未受圣洗}这一事实——他们却不是基督的同伴，而是信德的仇敌，律法的经师和法利塞人。由此可得出另一提示：既然主的\Emph{反对者}\Emph{拒绝}受圣洗，那么\Emph{跟随}主的人\Emph{都}受了圣洗，且不与自己的对手同心一意；尤其是，若有人依附何人，主曾藉见证高举若翰超乎那人之上，说：\Quote{凡妇人所生的，没有一个比洗者若翰更大。}\par

另有人提出一种相当牵强的说法：他们认为\Quote{宗徒们在小船中被海浪溅湿、淹没，便算受了圣洗；伯多禄自己在海上行走时也浸洗得够了。}然而，我认为，被海水的暴力溅湿或阻截是一回事，服从宗教规诫而受圣洗是另一回事。那只小船确实预表了教会：她在\Quote{海}中，即在世界中，被\Quote{波浪}，即迫害和诱惑骚扰；主如同在忍耐中沉睡，直到最后在圣人们的祈祷中被唤醒，祂便抑制世界，将宁静赐予祂自己的人。\par

现在，无论他们以何种方式受了圣洗，或至终未受圣洗——甚至主关于\Quote{一次沐浴}的说话，藉伯多禄的身份，仅是指\Emph{我们}而言——但要断定宗徒们的救恩是相当冒失的，因为\Emph{他们}享有首选的殊荣，以及此后不分彼此的亲密关系，或许能赋予他们简化的圣洗恩宠，因为他们（我以为）跟随那位惯于应许每个信徒救恩的主。祂会说：\Quote{你的\Emph{信德}救了你}，以及\Quote{你的罪得了赦免}，当然是在你相信之时，即使你还未\Emph{受圣洗}。若宗徒们缺少这一点，我不知道他们因何种事物之信，竟因主的一句话，\Emph{一个}人永远放弃了税关，\BibleRef{玛 9:9}\Emph{另一个}人抛弃了父亲、船和谋生的手艺，\BibleRef{玛 4:21}\Emph{又一个}人轻看了父亲的葬礼，在听到之前便遵行了主最崇高的诫命：\Quote{爱父亲或母亲超过我的，不配属于我。}\BibleRef{玛 10:37}\par

% structure: p0030 source=h2 target=section reason=relative-heading-rank; base=h2

\section{第十三章 另一异议：亚巴郎未受圣洗而悦乐天主。回应。旧事须让位于新事，圣洗如今成为律法}

于是，那些恶人挑起疑问。他们说：\Quote{圣洗对于信德已足够的人并非必须；因为亚巴郎并非藉水的圣事，而是藉信德悦乐了天主。}但在所有事情上，\Emph{后来}的具有决定性效力，\Emph{后续}的胜过在先的。诚然，在过去的年代，在主的苦难和复活之前，藉纯然的信德即可得救。但如今信德被扩大了，成为相信祂的诞生、苦难和复活的信德，于是圣事有了一项增补，即圣洗的封印行为；它在一定程度上包裹了先前裸露的信德，而这信德如今离开其固有律法便不能存在。因为施圣洗的\Emph{律法}已被\Emph{颁布}，仪式也被规定：\Quote{你们去}，\Emph{祂}说，\Quote{教导万民，因父、及子、及圣神之名给他们授圣洗。}将这定义\Quote{人除非由水和圣神重生，不能进入天国}与这律法相比，便将信德与圣洗的必然性联系起来了。因此，此后所有\Emph{成为}信徒的人都受了圣洗。\Emph{那时}，保禄相信后也受了圣洗；这就是主在使他遭受失明的惩罚时赐给他的诫命的意义，说：\Quote{起来，进入大马士革；在那里将指示你当做之事}，即——受圣洗，这是他所缺的唯一之事。除此一点外，他早已充分\Emph{学习并相信}\Quote{纳匝肋人}是\Quote{主、天主子}。\par

% structure: p0032 source=h2 target=section reason=relative-heading-rank; base=h2

\section{第十四章 论保禄声称他未被派遣去施圣洗}

但是，他们从\Emph{那位}宗徒本人那里提出一个\Emph{异议}，因为他曾说：\BibleQuote{因为基督差遣我，不是为施洗；}\BibleRef{格前 1:17} 仿佛由此论证圣洗就被废除了一般！因为\Emph{若是如此}，他为何给\Person{加约}、\Person{克黎斯颇}和\Person{斯德法纳}全家施洗呢？然而，即使基督没有差遣\Emph{他}施洗，祂却给了\Emph{其他}宗徒施洗的诫命。但这话是写给\Place{格林多}人，是针对那特定时期的情况；因为在他们中间发生了分裂和纷争，有人把\Emph{一切}都归于\Person{保禄}，有人归于\Person{阿颇罗}。为此，那位\Quote{缔造和平的}宗徒，唯恐他显得为自己索取一切\Emph{恩赐}，就说他被差遣\Quote{不是为施洗，而是为宣讲}。因为宣讲在先，施洗在后。所以，宣讲是\Emph{第一位的}：但我认为施洗对于那有权宣讲的人也是\Emph{合法的}。\par

% structure: p0034 source=h2 target=section reason=relative-heading-rank; base=h2

\section{第十五章 圣洗的统一性。关于异端和犹太人的圣洗的论述}

我不知道是否还有任何问题被提出来使圣洗成为争议。请允许我忆及以上遗漏的内容，以免我似乎打断了即将浮现的思路的中段。我们只有一位天主、一次圣洗、一个在天上的教会，这是按主的福音和宗徒的书信所说的，因为\Emph{他说}：\Quote{一个天主，一个圣洗，一个在天上的教会。}但必须承认，这个问题\Quote{关于异端者应遵守什么规则？}是值得讨论的。因为那断言是对\Emph{我们}说的。但异端者与我们的规范毫无关系，他们被绝罚的事实本身证明他们是外人。我没有义务在\Emph{他们}身上承认那对我有约束力的事，因为他们和我们没有同一个天主，也没有同一个基督。因此，他们的圣洗与我们的也不是同一个，因为不是同一个；他们既然没有正当的圣洗，无疑就根本没有圣洗；那不能\Emph{算为}有的东西，就是没有\Emph{有}。\BibleRef{训 1:15} 因此，他们也不能\Emph{领受}它，因为他们\Emph{没有它}。但这问题我们从希腊文中已有了更充分的讨论。我们只\Emph{一次}进入圣洗：\Emph{一次}罪被洗净，因为它们不应再重复。但犹太人的以色列每天沐浴，因为他每天被玷污；唯恐\Emph{污秽}也在\Emph{我们}中间发生，所以才制定了关于一次沐浴的规则。幸福的水，它\Emph{一次}洗净；它不戏弄罪人（以虚假的希望）；它不因重复污秽的感染而再次玷污它所洗净的人！\par

% structure: p0036 source=h2 target=section reason=relative-heading-rank; base=h2

\section{第十六章 第二次圣洗——以血}

我们确实也有\Emph{第二次}水泉（其本身与\Emph{第一次为一}），即\Emph{血}之水泉；关于此，主已经受洗后说：\Quote{我应当受一种洗礼，}当祂已经受过洗。因为祂\BibleQuote{以水及血而来，}\BibleRef{若一 5:6}正如\Person{若望}所记载；祂好藉水受洗，藉血得荣耀；为使我们同样\Emph{蒙召}藉\Emph{水}，\Emph{被拣选}藉\Emph{血}。这两个圣洗祂从祂被刺透的肋旁流出，使得那些相信祂血的人得以用水洗净；那些已用水洗净的人也得以同样饮血。这就是那圣洗，既在未领受水洗时代替水洗，又在其丧失后恢复它。\par

% structure: p0038 source=h2 target=section reason=relative-heading-rank; base=h2

\section{第十七章 关于赋予圣洗之权柄}

为了结束我们简短的论题，尚需提醒你们关于施洗和领洗的适当遵守。施洗的权利属于首席司祭（即主教）；其次，司铎和执事，但不得未经主教授权，以维护教会的尊荣，尊荣得以保存，平安亦得保存。此外，平信徒也有权；因为同样领受的，也可以同样给予。除非主教、司铎或执事在场，否则\Emph{其他}门徒被召叫\Emph{即从事这工作}。主的话不应被任何人隐藏；同样，圣洗作为天主之财产，可由所有人施行。但平信徒更应遵守恭敬和谦逊的规则——因为看见这些\Emph{权能}属于他们的上级——以免他们擅自承担主教的\Emph{特定}职务！效仿主教职务是分裂之母。最神圣的宗徒曾说过：\Quote{一切事都\Emph{许可}，但不都\Emph{有益}。}可以肯定，在\Emph{必要}的情况下，你利用（那规则）就够了，如果任何时候地点、时间或人的处境迫使你（这样做）；因为\Emph{那时}帮助者的坚定勇敢，当危难者的情况紧急时，是格外可接受的；因为他若拒绝给予他有权自由给予的东西，他将对一个人的丧亡有罪。但那轻率的女人，篡夺了教导的权柄，当然不会因此为自己生出施洗的权利，除非像从前一样出现某种新的野兽；以致正如一个废止了圣洗，另一个则凭借自己的权利施行它！但如果那些错误地归于保禄名下的著作，引用\Person{特克拉}的例证作为妇女教导和施洗的许可，让他们知道，在\Place{亚细亚}，那编写著作的司铎，似乎从自己的库存增加保禄的名声，被定罪并承认他出于对保禄的爱而这样做后，被革了职。因为那位连\Emph{女人}都未允许以过分大胆\Emph{学习}的人，怎么可能给\Emph{女性}\Emph{教导}和\Emph{施洗}的权柄呢！\BibleQuote{她们该沉默，}他说，\BibleQuote{并回家问自己的丈夫。}\BibleRef{格前 14:34}\par

% structure: p0040 source=h2 target=section reason=relative-heading-rank; base=h2

\section{第十八章 论应当施行圣洗的对象与时间}

然而，那些负有职责的人知道，圣洗不可轻率施行。\Quote{凡求你的，就给他}这句话有其自身所指，特别与施舍有关。相反地，这个\Emph{诫命}更应审慎思考：\BibleQuote{不要把圣物给狗，也不要把你们的珍珠丢在猪前}；\BibleRef{玛 7:6}以及\Quote{不可轻易给\Emph{任何人}覆手；不要分担别人的罪过。}如果斐理伯如此\Quote{轻易}地给那位太监施行了圣洗，让我们反思：已有明显而昭著的证据介入，表明主认为他配得。\BibleRef{宗 8:26}圣神曾命令斐理伯前往那条路；那位太监自己也不是闲散之辈，也非突然被急切的受洗愿望所攫取；而是在上圣殿祈祷之后，专心研读\WorkTitle{圣经}，就这样被恰当地发现——天主已不请自来地给他派出了一位宗徒，而这位宗徒又是圣神吩咐他去贴近太监的马车。他\Emph{正读的}经文恰好与他的信德吻合；\Emph{斐理伯}应请求被带上与他同坐；主被指明；信德没有迟延；水无需等候；圣事完成，宗徒被提走。\Quote{事实上，保禄也是\InnerQuote{迅速}领受了圣洗：}因为他的主人西满迅速认出他是\Quote{一个指定的蒙选之器}。天主的认可会在它之前发出确定的预兆；每一个\Quote{请求}都可能欺骗人，也可能受人欺骗。因此，根据每个人的情况、性格甚至年龄，延缓圣洗更为可取；然而，尤其在小孩子的身上。因为，如果（圣洗本身）并非如此必要，为何必要将代父母也推入危险？他们自己因必死性可能无法履行承诺，并因\Emph{他们所担保的人}邪恶倾向的发展而失望？主确实说过：\Quote{不要阻止他们到我这里来。}那么，让他们在成长时\Quote{来}，让他们在学习时\Quote{来}，在他们学习要去哪里时；让他们在能够认识基督时成为基督徒。为何纯洁的生命时期匆忙奔向\Quote{得罪赦}？在世俗事务上会更加谨慎：以致一个不被托付世俗财物的人，却被托付神圣之物！让他们知道如何\Quote{求取}救恩，这样你（至少）似乎给了\Quote{求者}。出于同样原因，未婚者也应延缓——在他们身上，\Emph{诱惑的土壤}已经预备好，无论是那些因成熟而\Emph{从未}结婚的人，还是那些因自由而\Emph{守寡}的人——直到他们结婚，或更坚强地持守节欲。如果有人理解圣洗的重大意义，他们会害怕领受它而非延缓它：健全的信德确保救恩。\par

% structure: p0042 source=h2 target=section reason=relative-heading-rank; base=h2

\section{第十九章 论最适合施行圣洗的时间}

逾越节提供比平时更庄严的圣洗日子；那时，主的苦难——我们藉以受洗——得以完成。以象征方式解释\Emph{这个事实}也不矛盾：当主要举行最后逾越节时，祂对被打发去准备的门徒们说：\Quote{你们会遇见一个提着水的人。}祂以\Emph{水}的记号指出举行逾越节的地点。之后，五旬节是施行圣洗最喜乐的时期；在此期间，主的复活也在门徒们中多次得到证实，主的来临之望也被间接指明，因为在那时，当祂被接回天上时，天使们告诉宗徒们：\Quote{祂要这样来，正如祂升天时一样}；当然，是在五旬节。然而，此外，当耶肋米亚说\Quote{我要在节期之日从地极聚集他们}时，他指的是逾越节和五旬节的日子，这确实是一个\Quote{节期之日}。然而，\Emph{每一天}都是主的日子；每一时辰，每一时刻，都适宜圣洗：如果\Emph{庄严性}有所不同，则\Emph{恩宠}并无区别。\par

% structure: p0044 source=h2 target=section reason=relative-heading-rank; base=h2

\section{第二十章 论领受圣洗的准备与领受后的行为}

凡将领受圣洗的人，当以多次祈祷、斋戒、屈膝、彻夜守更，并告解一切以往的罪过来祈求，好能表明甚至若翰洗礼的\Emph{含义}。\BibleQuote{他们受洗时，承认自己的罪。}我们若\Emph{现在}公开承认自己的不义或卑污，理当感恩；因为我们同时藉着克制肉身与神魂，为我们先前的罪做补赎，并为随后将至的诱惑预置防御的基础。主说：\BibleQuote{你们应当醒寤祈祷，免陷于诱惑。}\BibleRef{玛 26:41}我相信他们之所以陷于诱惑，是因为睡着了；所以他们在主被捕时离弃了祂，那位继续站在祂身边并动刀的，甚至三次否认了祂：因为话已先说了：\Quote{\Emph{未}经诱惑的人，不能进入天国。}主自己在\Emph{领洗}后立刻被诱惑包围，当时祂守了四十天斋戒。有人会说：\Quote{那么，\Emph{我们}也该在领洗\Emph{后}守斋。}然而，除非是喜乐的需要和救恩的感恩，谁禁止你呢？但以我浅薄的理解，主以比喻的方式将以色列人\Emph{对主所加的辱骂}回击给他们。因为百姓过海后，在旷野中漂泊了四十年，虽然在那里得到天主的供应，却更念及肚腹与咽喉，胜过念及天主。因此，主在受洗后退到旷野，以四十天的斋戒表明：天主的人生活\BibleQuote{不靠饼独生，而靠天主的圣言}；\BibleRef{玛 4:1}因饱食或食欲无度而生的诱惑，为禁欲所打破。因此，有福的\Emph{人们}，天主的恩宠正等待你们，当你们从这最神圣的重生之泉上来，首次在你们母亲家中与弟兄们举手时，向圣父祈求，向主祈求，使祂恩宠的殊恩\Emph{与}神恩的分施赐予你们。\BibleRef{格前 12:4}祂说：\BibleQuote{你们求，就必得到。}你们\Emph{已经}求了，也\Emph{已经}得了；你们\Emph{已经}敲门，门\Emph{已经}开了。惟愿你们祈求时，也记得罪人戴尔都良。\par

\SourceRecord{Translated by S. Thelwall.；Ante-Nicene Fathers；Vol. 3.；Edited by Alexander Roberts, James Donaldson, and A. Cleveland Coxe.；Buffalo, NY: Christian Literature Publishing Co.,；1885.；<http://www.newadvent.org/fathers/0321.htm>.}
